\section*{Hoofdstuk \ref{ch:b3:cancer-biology} --- Kankerbiologie}

\begin{solution}{exo:b3:cancer-biology:1}
Een \omterm{def:b3:cancer-biology:genes}{oncogen} is een functiewinstvorm van een gen dat de deling of de
overleving bevordert; één mutant allel volstaat (dominant):
\emph{RAS} (puntmutatie), \emph{MYC} (translocatie, vermenigvuldiging),
\emph{HER2}, \emph{BCR--ABL}. Een tumorsuppressor beteugelt de deling
of dwingt stilstand, dood of herstel af; beide allelen moeten verloren
gaan (recessief op celniveau): \emph{RB}, \emph{TP53}, \emph{APC},
\emph{BRCA1}.
\end{solution}

\begin{solution}{exo:b3:cancer-biology:2}
Aanhoudende deling: \emph{RAS}, \emph{MYC}. Ongevoeligheid voor
groeionderdrukkers: \emph{RB}, \emph{CDKN2A} (p16). De dood weerstaan:
\emph{BCL2}, \emph{TP53}. Onbeperkte deelbaarheid: de promotor van het
\omterm{def:b3:cancer-biology:telomerase}{telomerase}. Angiogenese: \emph{VEGF} aangezet via HIF (verlies van
\emph{VHL}). \omterm{def:b3:cancer-biology:metastasis}{Invasie} en \omterm{def:b3:cancer-biology:hallmarks}{metastase}: verlies van E-cadherine, de
epitheliaal-mesenchymale factoren. Genoominstabiliteit: de genen van
het \omterm{prop:b3:dna-repair:fidelity}{mismatchherstel}, \emph{\omterm{def:b3:dna-repair:dsb}{BRCA}}. Ontregelde stofwisseling:
\emph{MYC}, de \omterm{prop:b3:cancer-biology:pathways}{PI3K-route}. \omterm{prop:b3:cancer-biology:immune}{Immuunontwijking}: \emph{PD-L1}, verlies van
MHC. Ontsteking die de \omterm{def:b3:cancer-biology:hallmarks}{tumor} bevordert: de signalering van
\emph{NF-$\kappa$B}.
\end{solution}

\begin{solution}{exo:b3:cancer-biology:3}
Kinderen met een familiegeschiedenis kregen vroeg verscheidene \omterm{def:b3:cancer-biology:hallmarks}{tumoren}
in beide ogen; sporadische gevallen kregen later één \omterm{def:b3:cancer-biology:hallmarks}{tumor}. Hij leidde
daaruit af dat er in een cel twee gebeurtenissen nodig waren: de
familiaire patiënten hadden er één geërfd en hadden alleen een tweede,
somatische gebeurtenis nodig (vaak genoeg om in verscheidene cellen van
het netvlies op te treden), terwijl bij sporadische patiënten beide
somatisch in één cel moesten optreden --- zeldzaam, laat, enkelvoudig.
Dus een gen waarvan beide kopieën verloren moeten gaan: de eerste
tumorsuppressor.
\end{solution}

\begin{solution}{exo:b3:cancer-biology:4}
Een mutant van \emph{Salmonella} die histidine nodig heeft wordt met de
te toetsen stof zonder histidine uitgebracht; er ontstaan alleen
kolonies uit cellen waarin een nieuwe mutatie het gebrek heeft
hersteld, zodat hun aantal de mutageniteit meet. Er wordt leverextract
toegevoegd omdat veel carcinogenen onschadelijk zijn tot leverenzymen
ze in reactieve metabolieten omzetten, zoals in het lichaam gebeurt; de
bacteriën hebben die enzymen niet.
\end{solution}

\begin{solution}{exo:b3:cancer-biology:5}
De incidentie stijgt $100$-voudig terwijl de leeftijd verdubbelt:
$2^{k-1} = 100$, $k - 1 =
\log_{2}100 = 6.6$, dus ongeveer zeven of acht snelheidsbepalende
gebeurtenissen bij de naïeve lezing --- minder als klonale uitbreiding
tussen de treffers wordt toegelaten.
\end{solution}

\begin{solution}{exo:b3:cancer-biology:6}
$L = \sqrt{2Dc_{0}/q}$: voor $q = 0.01$ geeft dat $\sqrt{2\times 2\times 10^{-9}
\times 0.05/0.01} = \qty{141}{\micro m}$; voor $q = 0.03$
\qty{82}{\micro m}. Een snel groeiende \omterm{def:b3:cancer-biology:hallmarks}{tumor} ademt sneller, put de
zuurstof binnen een kortere afstand uit, en zijn midden sterft bij een
geringere omvang.
\end{solution}

\begin{solution}{exo:b3:cancer-biology:7}
$10^{11}\times 10^{-3n} < 1$ vergt $n > 3.7$: vier kuren. Na twee kuren
blijven er $10^{5}$ cellen over --- een tiende van een kubieke
millimeter, ver onder de $10^{9}$ die aantoonbaar zijn --- zodat de
\omterm{def:b3:cancer-biology:hallmarks}{tumor} volgens elke toets ``weg'' is terwijl er honderdduizend cellen
over zijn; de behandeling moet het geplande aantal kuren doorgaan en
niet tot de \omterm{def:b3:cancer-biology:hallmarks}{tumor} verdwijnt.
\end{solution}

\begin{solution}{exo:b3:cancer-biology:8}
Beide beschadigingen zetten dezelfde route aan (receptor $\to$ \omterm{def:b3:cancer-biology:genes}{Ras}
$\to$ ERK); zodra de ene aanwezig is geeft de andere geen verder
voordeel en wordt zij niet geselecteerd. Een EGFR-remmer blokkeert de
receptor; een cel die stroomafwaarts een KRAS-mutatie verwerft herstelt
het signaal terwijl de receptor geblokkeerd blijft, en wordt tijdens de
behandeling geselecteerd --- resistentie door omweg.
\end{solution}

\begin{solution}{exo:b3:cancer-biology:9}
E7 maakt E2F vrij van Rb en drijft de besmette epitheelcel zonder
groeifactoren door het \omterm{def:b3:cell-cycle-apoptosis:checkpoints}{restrictiepunt}; E6 vernietigt \omterm{def:b3:cancer-biology:genes}{p53}, zodat de
misplaatste deling en de schade die zij veroorzaakt noch stilstand noch
\omterm{def:b3:cell-cycle-apoptosis:apoptosis}{apoptose} uitlokken. Twee kenmerken worden in één keer geleverd, maar de
overige --- onsterfelijkheid, \omterm{def:b3:cancer-biology:metastasis}{invasie}, instabiliteit --- vergen verdere
somatische mutaties, die er decennia over doen om zich in de blijvend
besmette cellen op te hopen. Het vaccin wekt neutraliserende
antilichamen tegen het capside op en voorkomt de besmetting; wordt het
vóór de blootstelling gegeven, dan is er nooit een besmette cel om te
transformeren.
\end{solution}

\begin{solution}{exo:b3:cancer-biology:10}
Na de $i$de treffer groeit de kloon met $g$, zodat er $g^{i}$ keer
zoveel cellen de volgende treffer lopen; de kans op de volledige reeks
in de lijn die van één oorspronkelijke cel afstamt is
$g\cdot g^{2}\cdots$ --- bij constante $g$ per treffer wordt het
cumulatieve risico $N(ut)^{k}
g^{k-1}$ op volgordeafhankelijke factoren na, en houdt de incidentie de
macht $t^{k-1}$ met een voorfactor die $g^{k-1}$ groter is. Groeien de
uitbreidingen zelf met de tijd (exponentieel groeiende klonen), dan
wordt de curve nog steiler, zodat een helling van $6$ door minder dan
zeven aandrijvers wordt voortgebracht: de $k$ van Armitage en Doll is
een bovengrens voor het aantal snelheidsbepalende aandrijvers, en
gesequencede \omterm{def:b3:cancer-biology:hallmarks}{tumoren} tonen er twee tot acht.
\end{solution}

\begin{solution}{exo:b3:cancer-biology:11}
(a) Middel A alleen staat tegenover $\mu_{A}N = 100$ reeds aanwezige
resistente cellen: $P_{0} = e^{-100} \approx 0$; de A-resistente kloon
overleeft en groeit aan tijdens de achttien weken A, en tegen de tijd
dat B begint is zij weer een grote populatie, zodat ook B tegenover
$\mu_{B}N' \gg 1$ staat --- de reeks faalt tweemaal. (b) Samen: een
dubbel resistente cel ontstaat met $10^{-14}$ per deling, dus
$\mu_{A}\mu_{B}N = 10^{-5}$ en $P_{0} =
0.99999$. Vanaf het begin combineren, wanneer $N$ het kleinst is, is de
regel omdat juist het product van twee kleine snelheden de resistentie
onwaarschijnlijk maakt.
\end{solution}

\begin{solution}{exo:b3:cancer-biology:12}
Met $100$ keer meer cellen en dezelfde $u$ en $k$ zou het cumulatieve
risico $N(ut)^{k}$ $100$ keer hoger zijn --- elke olifant zou jong aan
\omterm{def:b3:cancer-biology:hallmarks}{kanker} moeten sterven. Dat gebeurt niet, dus moeten $u$ of $k$
verschillen: twintig kopieën van \emph{TP53} maken de apoptotische
reactie op schade sterker en halen mutante cellen weg voordat zij
uitbreiden (wat de effectieve $u$ voor die stap verlaagt); grote dieren
hebben ook lagere mutatiesnelheden per cel (beter herstel, minder
delingen per cel per jaar omdat hun cellen zich trager vernieuwen), en
hebben mogelijk meer treffers nodig (een extra laag suppressoren). De
onderdrukking van \omterm{def:b3:cancer-biology:hallmarks}{kanker} is met de lichaamsgrootte meegeëvolueerd.
\end{solution}

\begin{solution}{pb:b3:cancer-biology:1}
\textbf{1.} $10^{12}/10^{3} = 10^{9}$ cellen; $\log_{2}10^{9} \approx 30$
verdubbelingen.
\textbf{2.} $30\times 100 = 3000$ dagen, ongeveer $8$ jaar.
\textbf{3.} \qty{1}{kg} $\approx 10^{12}$ cellen: tien verdubbelingen
meer, $1000$ dagen, minder dan drie jaar --- een derde van de tijd tot
de ontdekking.
\textbf{4.} De verdubbelingstijd moet in het begin ongeveer $60$ dagen
zijn geweest, of de \omterm{def:b3:cancer-biology:hallmarks}{tumor} groeide sneller toen hij klein was en
vertraagde naarmate hij groeide (het gewone patroon, naarmate aanvoer
en ruimte tekortschieten).
\textbf{5.} Een cyclus van $2$ dagen zou de populatie elke $2$ dagen
verdubbelen; een verdubbelingstijd van $100$ betekent dat de
nettogroei $2\,\%$ van de geboortesnelheid is: vrijwel elke cel die
geboren wordt staat tegenover een die sterft, of er delen er maar
weinig --- geboorte en dood bijna in evenwicht, met de groei als klein
verschil.
\textbf{6.} De middelen doden delende cellen, dus verliest een \omterm{def:b3:cancer-biology:hallmarks}{tumor} met
een klein deel in cyclus per kuur een kleiner deel; maar de snelle
\omterm{def:b3:cancer-biology:hallmarks}{tumor}, die meer heeft verloren, groeit tussen de kuren ook weer aan.
Per saldo zijn het de snelle \omterm{def:b3:cancer-biology:hallmarks}{tumoren} (leukemieën, lymfomen,
teelbalkanker) die de chemotherapie geneest.
\textbf{7.} Verhouding $300$ bij leeftijdsverhouding $80/30 = 2.67$: $k - 1 =
\ln 300/\ln 2.67 = 5.8$, $k \approx 7$.
\textbf{8.} $N(ut)^{k} = 10^{8}\times (7\times 10^{-5})^{7} \approx
10^{-21}$: onzinnig tegenover een levenslang risico van bijna $1$ op
$3$. Wat ontbreekt: de klonale uitbreiding die het aantal cellen dat de
treffer loopt na elke stap vermenigvuldigt, mutatiesnelheden die door
instabiliteit zijn verhoogd, veel meer delingen (stamcellen die
decennia doordelen), en minder echte aandrijvers.
\textbf{9.} Vermenigvuldigen met $g^{k-1} = 100^{6} = 10^{12}$ geeft
$10^{-9}$ --- met zeven gebeurtenissen nog altijd minuscuul; met vier
gebeurtenissen en uitbreidingen geeft
$10^{8}\times (7\times 10^{-5})^{4}\times 10^{6} \approx 2\times
10^{-3}$, de juiste orde. Het ware beeld is weinig aandrijvers en grote
uitbreidingen.
\textbf{10.} $(ut)^{k-1}/(ut)^{k} = 1/(ut) = 1/(10^{-6}\times 40) =
2.5\times 10^{4}$: de incidentie van een draagster op veertigjarige
leeftijd is tienduizenden keren die van het sporadische geval ---
erfelijke \omterm{def:b3:cancer-biology:hallmarks}{kankers} slaan jong toe.
\textbf{11.} Eén gebeurtenis tienmaal sneller: incidentie $\times 10$;
twee: $\times 100$. Het waargenomen risico van vijftien tot
vijfentwintig keer wijst erop dat de rook op meer dan één stap
inwerkt.
\textbf{12.} Bij het stoppen valt $u$ terug voor de gebeurtenissen die
nog moeten komen, zodat de incidentie niet meer zo steil klimt; maar de
treffers die zich al in de cellen hebben opgehoopt blijven, zodat de
cellen van de ex-roker dichter bij \omterm{def:b3:cancer-biology:hallmarks}{kanker} staan dan die van iemand die
nooit heeft gerookt en het extra risico blijft, bevroren in plaats van
uitgewist.
\textbf{13.} $L = \sqrt{2\times 2\times 10^{-9}\times 0.05/0.03} =
\qty{82}{\micro m}$.
\textbf{14.} Volledig van zuurstof voorzien tot een doorsnede $2L \approx
\qty{160}{\micro m}$. Een bolletje van \qty{2}{mm} heeft een levende
rand van \qty{82}{\micro m}: het levende deel is
$1 - (918/1000)^{3} = 0.23$.
\textbf{15.} De cellen halverwege liggen \qty{75}{\micro m} van een vat,
net binnen $L$ --- in beginsel van zuurstof voorzien, maar tumorvaten
voeren minder zuurstof aan en stromen met horten en stoten, zodat de
cellen halverwege hypoxisch zijn. Hypoxische cellen zijn twee tot drie
keer beter bestand tegen straling (zuurstof is nodig om de schade
blijvend te maken), en de hypoxie lokt p53-afhankelijke \omterm{def:b3:cell-cycle-apoptosis:apoptosis}{apoptose} uit,
zodat zij de cellen selecteert die \omterm{def:b3:cancer-biology:genes}{p53} hebben verloren.
\textbf{16.} De dichtheid halveren spreidt de vaten met $\sqrt{2}$: zij
liggen \qty{212}{\micro m} uit elkaar, de cellen halverwege
\qty{106}{\micro m} weg, voorbij $L$: zij sterven of worden, als zij
hypoxisch overleven, agressiever en angiogener.
\textbf{17.} $30/2 = 15$ keer meer glucose. \omterm{def:b3:cancer-biology:hallmarks}{Tumoren} concentreren daarom
glucose-analogen, en een scan met positronenuitzendend
fluordeoxyglucose laat hen oplichten.
\textbf{18.} Uitgehongerde \omterm{def:b3:cancer-biology:hallmarks}{tumoren} krimpen tot de door diffusie
begrensde omvang en blijven sluimerend en hypoxisch bestaan, of nemen
bestaande vaten in gebruik; de hypoxie selecteert agressieve klonen.
Maar het snoeien van het chaotische vaatstelsel van de \omterm{def:b3:cancer-biology:hallmarks}{tumor} verlaagt de
druk in het weefsel en maakt de stroom gelijkmatiger, zodat de
chemotherapie beter doordringt --- de combinatie werkt waar geen van
beide dat alleen doet.
\textbf{19.} $10^{9}\times 10^{-2n} < 1$: $n > 4.5$, vijf kuren.
\textbf{20.} $21$ dagen $= 0.21$ verdubbeling, aangroei $\times 1.16$;
nettofactor per cyclus $0.0116$; $\ln 10^{-9}/\ln 0.0116 = 4.65$: nog
steeds vijf cycli.
\textbf{21.} $10^{-7}\times 10^{9} = 100$ resistente cellen te
verwachten; $P_{0} = e^{-100} \approx 0$: de resistentie is zeker.
\textbf{22.} $\mu_{1}\mu_{2}N = 10^{-14}\times 10^{9} = 10^{-5}$;
$P_{0} = 0.99999$. Bij $10^{12}$ cellen $10^{-2}$, dus $P_{0} = 0.99$.
\textbf{23.} $10^{6}$ cellen: één middel $\mu N = 0.1$, $P_{0} = 0.90$;
twee middelen $10^{-8}$, $P_{0} \approx 1$. Behandel wanneer de
populatie klein is, en behandel met combinaties.
\textbf{24.} De T-cellen vallen vele neoantigenen tegelijk aan, zodat
geen enkele mutatie hun ontsnapt; ontsnappen vergt het verlies van de
antigeenpresentatie zelf (MHC, $\beta_{2}$-microglobuline) of het
onderdrukken van de T-cellen, een andere en zeldzamere soort
gebeurtenis. Dezelfde hoge mutatielast die een \omterm{def:b3:cancer-biology:hallmarks}{tumor} zeker resistent
maakt tegen een kinaseremmer geeft hem vele neoantigenen: \omterm{def:b3:cancer-biology:hallmarks}{tumoren}
zonder \omterm{prop:b3:dna-repair:fidelity}{mismatchherstel} en \omterm{def:b3:cancer-biology:hallmarks}{tumoren} die door ultraviolet licht of rook
zijn opgewekt reageren het best.
\textbf{25.} Ongeveer $8$ jaar tot de ontdekking; levende \omterm{def:b3:cancer-biology:hallmarks}{tumor} tot
\qty{82}{\micro m} diep rond een vat; $P_{0} = 0.99999$ dat een
combinatie van twee middelen in een \omterm{def:b3:cancer-biology:hallmarks}{tumor} van \qty{1}{cm^3} geen
resistente cel tegenkomt.
\end{solution}
